% ch01  §1 引言 / §2.1–§2.3（原文 585–962 行）

\section{引言}

shtuka 这一概念由 Drinfeld \cite{Dr-FBun} 以「$F$-丛」之名引入，其中 $F$ 指 Frobenius。事实上，shtuka 就是带某种 Frobenius 结构的丛。「shtuka」一词源自俄语，意为「小玩意」或「东西」。shtuka 推广了 Drinfeld 更早引入的「椭圆模」概念，后者是椭圆曲线的类比物。一般的 shtuka 可以看成动机的一种版本：它不定义在特征 $p$ 之中，但其系数域的特征为 $p$。

shtuka 的模空间或模叠是 Shimura 簇在函数域上的类比物。因此，它们在建立函数域上 Langlands 对应的各种情形时一直起着关键作用。

Drinfeld \cite{Dr-Pet, Dr-comp} 用秩 $2$ shtuka 的模空间证明了整体函数域上 $\GL_{2}$ 的 Ramanujan–Petersson 猜想与 Langlands 猜想。

Laumon、Rapoport 与 Shtuhler \cite{LRS} 用除代数上的 shtuka 证明了函数域上除代数的整体 Langlands 猜想的一部分，并证明了局部函数域上 $\GL_{n}$（任意 $n$）的局部 Langlands 猜想。

L. Lafforgue \cite{LLaff} 用秩 $n$ shtuka 的模叠及其紧化，对任意 $n\in\NN$ 证明了整体函数域上 $\GL_{n}$ 的 Ramanujan–Petersson 猜想与 Langlands 猜想。

Varshavsky \cite{Var} 把 shtuka 的概念从向量丛的情形推广到主 $G$-丛的情形，其中 $G$ 是基域 $\FF_{q}$ 上任一分裂约化群。我们称这样的对象为 $G$-shtuka。他证明了关于 $G$-shtuka 模叠的几何的基本结果。

更晚近地，V. Lafforgue \cite{VLaff} 利用 Varshavsky 的 $G$-shtuka 模空间，结合若干巧妙的想法（其中最重要的是游走算子），证明了函数域上任一连通约化群 $G$ 的整体 Langlands 猜想的「自守到 Galois」方向。

局部版的 shtuka 虽不在本文讨论范围之内，却在证明局部 Langlands 猜想时起重要作用。在 \cite{GL} 中，Genestier 与 V. Lafforgue 用局部 shtuka 证明了局部函数域上任一连通约化群 $G$ 的局部 Langlands 猜想的「表示到 Galois」方向。另一项令人兴奋的进展是 Scholze 关于混合特征局部域的局部 shtuka 理论，其模空间推广了 Rapoport–Zink 空间，参见 \cite{S-Berkeley}。

本文是一篇综述。我们将在 \S\ref{s:def} 中从头介绍 $\FF_{q}$ 上分裂约化群的 shtuka 概念；接着在 \S\ref{s:sib} 中描述它与 Drinfeld 模的关系，以及它与 Deligne–Lusztig 簇、有限域上动机的类比；然后在 \S\ref{s:geom} 与 \S\ref{s:coho} 中研究 shtuka 模叠上的一些重要几何构造，例如 Hecke 对应与部分 Frobenius，并讨论它们对模叠上同调的影响。这些内容为 C. Xue \cite{Xue} 的讲座提供背景。最后一节 \S\ref{s:adv} 提及若干值得进一步研究的有趣课题，例如模叠的紧化，以及把 shtuka 模叠上的特殊圈与 $L$-函数、Eisenstein 级数的高阶导数联系起来的新近工作。

关于 shtuka 已有出色的综述文章，例如 \cite{Laumon} 与 \cite{NDT-Sht}，作者正是从中习得这套理论。研究论文 \cite{Var} 与 \cite{ARH} 可作为证明细节的详尽参考。最后，本文很大一部分遵循突破性论文 \cite{VLaff} 的处理方式，读者可从中找到更多细节及其在 Langlands 纲领中的应用。

%%% thanks %%%

\subsection*{致谢}

感谢 2022 年 IHES 暑期学校（Langlands 纲领专题）的组织者邀请我就此主题作报告。感谢 Jianqiao Xia 与审稿人仔细阅读草稿并提出有益的建议。



\section{基本定义}\label{s:def}

\subsection{曲线上的 $G$-主丛}
设 $k$ 为有限域 $\FF_{q}$。对两个 $k$-叠 $\cX$ 与 $\cY$，用 $\cX\times\cY$ 表示它们在 $\Spec k$ 上的纤维积。

设 $X$ 是 $k$ 上几何连通的光滑射影曲线，$F=k(X)$ 为其函数域，$|X|$ 为 $X$ 的闭点集。对 $x\in |X|$，记 $\cO_{x}$、$\fm_{x}$、$F_{x}$、$k_{x}$ 分别为 $X$ 在 $x$ 处的完备化局部环及其极大理想、分式域与剩余域。

\ssszh{$G$-主丛}
设 $G$ 是 $k$ 上分裂连通约化群。$k$-概形 $Y$ 上的一个 $G$-主丛，是指一个概形 $\cE$，带 $G$ 的右作用以及 $G$-不变的映射 $\pi: \cE\to Y$，使得在 $Y$ 上 étale 局部地，$\cE$ 与 $G\times Y$（带 $G$ 的右平移作用）$G$-等变同构。

令 $\Bun_{G}: \Sch_{k}\to \Gpd$（其中 $\Sch_{k}$ 是 $k$-概形范畴，$\Gpd$ 是群胚范畴）为如下函子：把 $k$-概形 $S$ 对应到 $X\times S:=X\times_{k}S$ 上 $G$-主丛的群胚，其态射为 $G$-主丛的同构。则 $\Bun_{G}$ 可由 $k$ 上局部有限型的光滑 Artin 叠表示。

\begin{exam} 设 $G=\GL_{n}$。概形 $Y$ 上 $\GL_{n}$-主丛的群胚与 $Y$ 上秩 $n$ 向量丛的群胚等价。因此 $\Bun_{G}$ 同构于 $X$ 上秩 $n$ 向量丛的模叠 $\Bun_{n}$。

若 $G=\SL_{n}$，则 $\Bun_{G}(S)$ 参数化偶 $(\cV, \l)$，其中 $\cV$ 是 $X\times S$ 上秩 $n$ 的向量丛，$\l$ 是同构 $\l: \cO_{X\times S}\isom \det(\cV):=\wedge^{n}(\cV)$。
\end{exam}



\begin{exam} 设 $G=\Sp_{2n}$。概形 $Y$ 上 $\Sp_{2n}$-主丛的群胚与 $Y$ 上辛向量丛 $(\cV, \om)$ 的群胚等价。这里 $\cV$ 是 $Y$ 上秩 $2n$ 的向量丛，$\om: \cV\ot \cV\to \cO_{Y}$ 是完美交错配对。因此 $\Bun_{G}$ 同构于 $X$ 上秩 $2n$ 辛向量丛的模叠。
\end{exam}



\ssszh{水平结构}\label{sss:level}

设 $D\subset X$ 是有限子概形，$S$ 是 $k$-概形。$X\times S$ 上 $G$-主丛 $\cE$ 的一个 $D$-水平结构，是指限制所得 $D\times S$ 上 $G$-主丛 $\cE|_{D\times S}$ 的一个平凡化。令 $\Bun_{G,D}: \Sch_{k}\to \Gpd$ 为如下函子：它在 $S$ 处的取值是偶 $(\cE,\io_{D})$ 的群胚，其中 $\cE$ 是 $X\times S$ 上的 $G$-主丛，$\io_{D}$ 是 $\cE$ 的 $D$-水平结构。则 $\Bun_{G,D}$ 可由 Artin 叠表示，且遗忘映射 $\Forg_{D}: \Bun_{G,D}\to \Bun_{G}$ 是一个 $G_{D}:=\Res_{D/k}(G\times D)$-主丛。


\begin{exam} 设 $G=\Sp_{2n}$，并对某个 $x\in |X|$ 取 $D=x$。对任意 $k$-概形 $S$，$\Bun_{G,x}(S)$ 是多元组 $(\cV, \om, \{e_{1},\cdots, e_{n}, f_{1},\cdots, f_{n}\})$ 的群胚，其中 $(\cV,\om)$ 是 $X\times S$ 上秩 $2n$ 的辛向量丛，$e_{i}$ 与 $f_{i}$ 是 $\cV|_{D\times S}$ 的截面，且对所有 $1\le i,j\le n$ 有 $\om(e_{i},e_{j})=0$、$\om(f_{i},f_{j})=0$、$\om(e_{i},f_{j})=\d_{ij}$（这些都是 $\cO_{D\times S}$ 的截面之间的等式）。
\end{exam}


考虑其它的水平结构也有用处。设 $x\in |X|$。记 $L^{+}_{x}G$ 为 $G$ 在 $x$ 处的形式弧群：它是 $k_{x}$ 上的群概形，其 $R$-点（$R$ 为 $k_{x}$-代数）为 $G(R\hotimes_{k_{x}} \cO_{x})$。记 $L_{x}G$ 为 $G$ 在 $x$ 处的形式环路群：它是 $k_{x}$ 上的群 ind-概形，其 $R$-点为 $G(R\hotimes_{k_{x}} F_{x})$。参见 \cite[\S1.a]{PR}。

对 $n\in \NN$，令 $D_{nx}=\Spec \cO_{x}/\fm^{n}_{x}$。我们有 $L^{+}_{x}G$ 的截断版本 $G_{nx}:=\Res_{D_{nx}/k_{x}}(G\times D_{nx})$。令 $r_{nx}: L^{+}_{x}G\to G_{nx}$ 为约化映射。称
\begin{equation*}
L^{(n)}_{x}G=\ker(r_{nx})
\end{equation*}
为 $L^{+}_{x}G$ 的第 $n$ 个主同余子群。

设 $\bK_{x}\subset L^{+}_{x}G$ 是 $k_{x}$ 上的子群概形，且具有形状 $r_{nx}^{-1}(\ov\bK_{x,n})$，其中 $n\in\NN$，$\ov\bK_{x,n}\subset G_{nx}$ 是 $k_{x}$-子群。定义 $\Bun_{G,\bK_{x}}$ 为商叠
\begin{equation*}
\Bun_{G,\bK_{x}}:=\Bun_{G, nx}/(\Res_{k_{x}/k}\ov\bK_{x,n}).
\end{equation*}
注意 $\Res_{k_{x}/k}G_{nx}=\Res_{D_{nx}/k}(G\times D_{nx})$ 作用在 $\Bun_{G, nx}$ 上，因此其子群 $\Res_{k_{x}/k}\ov\bK_{x,n}$ 也作用。容易看出，$\Bun_{G,\bK_{x}}$ 在相差典范同构的意义下与 $n$ 的选取无关。



\begin{exam} 设 $\bI_{x}$ 是 $L^{+}_{x}G$ 的 Iwahori 子群，即它是某个 Borel 子群 $B_{x}\subset G\ot k_{x}$ 在 $r_{x}: L^{+}_{x}G\to G\ot k_{x}$ 下的原像。则 $\Bun_{G,\bI_{x}}$ 的 $S$-点与偶 $(\cE, \cF_{x})$ 是同一份数据，其中 $\cE$ 是 $X\times S$ 上的 $G$-主丛，$\cF_{x}$ 是限制 $\cE|_{\{x\}\times S}$ 的一个 Borel 约化，也就是丛 $\cE|_{\{x\}\times S}\times^{G}\fl_{G}\to \{x\}\times S$ 的一个截面（$\fl_{G}$ 是 $G$ 的旗簇）。由后一描述可见，$\Bun_{G,\bI_{x}}$ 典范地不依赖于 Borel 子群 $B_{x}$ 的选取。
\end{exam}

更一般地，当 $\bK_{x}\subset L_{x}G$ 不必包含在 $L^{+}_{x}G$ 中，而是以 $L^{+}_{x}G$ 的某个同余子群为余维有限的正规子群时，仍可定义模叠 $\Bun_{G,\bK_{x}}$。这样的 $\bK_{x}$ 称为 $x$ 处的 {\em 水平群}（level group）。为构造 $\Bun_{G,\bK_{x}}$，令 $\Bun_{G, \infty x}$ 为塔 $(\Bun_{G,nx})_{n\in\NN}$ 的逆向极限。可以证明，$\Bun_{G,\infty x}$ 上的 $L^{+}_{x}G$-作用典范地延拓为 $L_{x}G$ 的作用。于是定义 $\Bun_{G,\bK_{x}}$ 为商 $\Bun_{G,\infty x}/(\Res_{k_{x}/k}\bK_{x})$。为看出它是 Artin 叠，取一个在 $\bK_{x}$ 中正规的主同余子群 $L^{(n)}_{x}G$，令 $\ov \bK_{x,n}=\bK_{x}/L^{(n)}_{x}G$，并把 $\Bun_{G,\infty x}/(\Res_{k_{x}/k}\bK_{x})$ 等同于 $\Bun_{G,nx}/(\Res_{k_{x}/k}\ov\bK_{x,n})$。特别地，上述构造可施于抛孔子群 $\bP_{x}\subset L_{x}G$。

\begin{rmk} 对一般的水平结构 $\bK_{x}$，习惯上把 $\Bun_{G,\bK_{x}}$ 的 $S$-点说成「$X\times S$ 上带 $x$ 处 $\bK_{x}$-水平结构的 $G$-主丛」。然而这一术语可能引起误解：当 $\bK_{x}$ 不必包含在 $L^{+}_{x}G$ 中时，$\Bun_{G,\bK_{x}}$ 的 $S$-点并不诱导整个 $X\times S$ 上的 $G$-主丛，而只诱导 $(X-\{x\})\times S$ 上的 $G$-主丛。下面的例说明了这一点。
\end{rmk}



\begin{exam}\label{ex:SL parahoric} 设 $G=\SL_{n}$，并固定 $x\in |X|$。在共轭意义下，$L_{x}G$ 有 $n$ 个极大抛孔子群 $\bP^{(i)}_{x}$，由 $i\in \ZZ/n\ZZ$ 参数化。若令 $\wt G=\GL_{n}$，则可取 $\bP^{(i)}_{x}=g_{i}(L^{+}_{x}\wt G)g_{i}^{-1}\cap L_{x}G$，其中 $g_{i}$ 是 $\wt G(F_{x})$ 中任一满足 $\val_{F_{x}}(\det(g_{i}))\equiv i\pmod{n}$ 的元素。

对 $\Bun_{G, \bP^{(i)}_{x}}$ 有如下解释。考虑把 $\cV$ 送到 $\wedge^{n}(\cV)$ 的映射 $\det: \Bun_{\wt G}\to \Pic_{X}$。则存在叠的典范同构
\begin{equation*}
\Bun_{G, \bP^{(i)}_{x}}\cong {\det}^{-1}(\cO(-ix)), \quad i\in \ZZ.
\end{equation*}
当然，与 $\cO(x)$ 作张量积会把 $\det^{-1}(\cO(-ix))$ 与 $\det^{-1}(\cO(-(i-n)x))$ 等同起来，故上式右端的同构类只依赖于 $i$ 在 $\ZZ/n\ZZ$ 中的像。

换言之，对 $i\in \ZZ$，$\Bun_{G, \bP^{(i)}_{x}}$ 的 $S$-点分类偶 $(\cV,\l)$，其中 $\cV$ 是 $X\times S$ 上秩 $n$ 的向量丛，$\l$ 是同构 $\cO(-ix)\isom \wedge^{n}(\cV)$。
\end{exam}


上述定义显然可推广到在不止一个、而是有限多个点处给水平结构的情形。若 $\Sig\subset |X|$ 是有限子集，$\bK:=(\bK_{x})_{x\in \Sigma}$ 是每个 $x\in\Sig$ 处水平群的族，则把所得的模叠记作 $\Bun_{G,\bK}$。



\subsection{$G$-主丛的 Hecke 叠}
以下固定有限个位点组成的集合 $\Sigma\subset |X|$，并对每个 $x\in \Sigma$ 固定一个水平群 $\bK_{x}$。令 $U=X-\Sigma$，用 $\bK$ 记水平群族 $(\bK_{x})_{x\in \Sigma}$。于是模叠 $\Bun_{G,\bK}$ 有定义。


\ssszh{Hecke 叠}\label{sss:Hk}

设 $I$ 是有限集。令 $\Hk^{I}_{G,\bK}: \Sch_{k}\to \Gpd$ 为如下函子：它在 $k$-概形 $S$ 处的取值是多元组 $\xi=((x_{i})_{i\in I}, \cE_{0},\cE_{1} ,\a)$ 的群胚，其中
\begin{enumerate}
\item 对每个 $i\in I$，$x_{i}: S\to U$ 是态射，其图像为 $\G(x_{i})\subset X\times S$。这些 $x_{i}$ 称为 $\xi$ 的 {\em 腿}（leg）。
\item $\cE_{0},\cE_{1}\in \Bun_{G,\bK}(S)$，即 $\cE_{0}$ 与 $\cE_{1}$ 是 $X\times S$ 上的 $G$-主丛，并沿 $\{x\}\times S$ 带有 $\bK_{x}$-水平结构。
\item $\a$ 是带 $\bK$-水平结构的 $G$-主丛之间的同构 $\cE_{0}|_{X\times S-\cup_{i\in I}\G(x_{i})}\isom \cE_{1}|_{X\times S-\cup_{i\in I}\G(x_{i})}$。
\end{enumerate}

按定义，$\Hk_{G,\bK}^{I}$ 带有映射
\begin{equation*}
\Hk^{I}_{G,\bK}\to U^{I}.
\end{equation*}

当 $I=\vn$ 时，$\Hk^{I}_{G,\bK}=\Bun_{G,\bK}$。当 $\Sig=\vn$（即 $\bK$ 无数据）时，把 $\Hk^{I}_{G,\bK}$ 简写作 $\Hk^{I}_{G}$。


\ssszh{Hecke 叠的迭代版本}\label{sss:it Hk}
设 $I=I_{1}\coprod \cdots\coprod I_{r}$ 是有限集 $I$ 的分解。把同构 $\a$ 细化为 $r$ 步，便得到 Hecke 叠 $\Hk^{I}_{G,\bK}$ 的变体 $\Hk^{(I_{1},\cdots, I_{r})}_{G,\bK}$。$\Hk^{(I_{1},\cdots, I_{r})}_{G,\bK}$ 的 $S$-点由多元组 $((x_{i})_{i\in I}, (\cE_{j})_{0\le j\le r} ,(\a_{j})_{1\le j\le r})$ 组成，其中
\begin{enumerate}
\item 对每个 $i\in I$，$x_{i}: S\to U$ 是态射，其图像为 $\G(x_{i})\subset X\times S$。
\item $\cE_{0},\cdots, \cE_{r}\in \Bun_{G,\bK}(S)$。
\item 对 $1\le j\le r$，$\a_{j}$ 是带 $\bK$-水平结构的 $G$-主丛之间的同构
$$\a_{j}: \cE_{j-1}|_{X\times S-\cup_{i\in I_{j}}\G(x_{i})}\isom \cE_{j}|_{X\times S-\cup_{i\in I_{j}}\G(x_{i})}.$$
\end{enumerate}

记录丛 $\cE_{j}$ 给出态射
\begin{equation*}
p_{j}: \Hk^{(I_{1},\cdots, I_{r})}_{G,\bK}\to \Bun_{G,\bK}, \quad j=0,1,\cdots, r.
\end{equation*}

忘掉中间丛 $\cE_{1},\cdots, \cE_{r-1}$，并取复合 $\a=\a_{r}\c\cdots \c\a_{1}$，便得到映射
\begin{equation*}
\pi^{(I_{1},\cdots, I_{r})}_{I}: \Hk^{(I_{1},\cdots, I_{r})}_{G,\bK}\to \Hk^{I}_{G,\bK}.
\end{equation*}
该映射在开子概形 $U^{(I_{1},\cdots, I_{r})}_{\disj}\subset U^{I}$ 上是同构；此开子概形由满足下述条件的 $(x_{i})_{i\in I}$ 组成：对属于不同部分 $I_{j}$ 与 $I_{j'}$ 的 $i$ 与 $i'$，$x_{i}$ 与 $x_{i'}$ 的图像不相交。

更一般地，若每个 $I_{j}$ 都有分解 $I_{j}=I_{j,1}\coprod \cdots\coprod I_{j,s_{j}}$，则 $I$ 有更细的分解 $I=I_{1,1}\coprod \cdots\coprod I_{1,s_{1}}\coprod I_{2,1}\coprod\cdots\coprod I_{2,s_{2}}\coprod \cdots\coprod I_{r,s_{r}}$。把这个加细记作 $I=I'_{1}\coprod \cdots\coprod I'_{s}$（故 $s=s_{1}+\cdots+s_{r}$），其中各部分的顺序与上面相同。我们得到映射
\begin{equation*}
\pi^{(I'_{1},\cdots, I'_{s})}_{(I_{1},\cdots, I_{r})}: \Hk^{(I'_{1},\cdots, I'_{s})}_{G,\bK}\to \Hk^{(I_{1},\cdots, I_{r})}_{G,\bK}
\end{equation*}
它只记录丛 $\cE_{0},\cE_{s_{1}},\cE_{s_{1}+s_{2}},\cdots, \cE_{s}$ 以及复合 $\a_{j}:=\a_{j,s_{j}}\c\cdots\c\a_{j,1}: \cE_{j-1}\dr\cE_{j}$，$1\le j\le r$。该映射在 $\prod_{j=1}^{r}U^{(I_{j,1},\cdots, I_{j,s_{j}})}_{\disj}\subset \prod_{j=1}^{r}U^{I_{j}}=U^{I}$ 上是同构。



\ssszh{仿射 Grassmann 簇}\label{sss:Gr}
关于这一主题的更多细节，读者可参阅 X. Zhu 的综述文章 \cite{Zhu-Satake}。$G$ 的仿射 Grassmann 簇 $\Gr_{G}$ 按定义是 $k$ 上的 ind-概形，它表示 $k$-代数上函子 $R\mapsto (LG)(R)/(L^{+}G)(R)=G(R\lr{t})/G(R\tl{t})$ 的 fpqc 层化。或者，若令 $D_{R}=\Spec R\tl{t}$、$D^{\times}_{R}=\Spec R\lr{t}$，则 $\Gr_{G}(R)$ 分类偶 $(\cE, \io)$，其中 $\cE$ 是 $D_{R}$ 上的 $G$-主丛，$\io$ 是 $\cE|_{D^{\times}_{R}}$ 的平凡化，即同构 $\io: \cE|_{D^{\times}_{R}}\isom \cE_{\triv}|_{D^{\times}_{R}}$，这里 $\cE_{\triv}$ 是 $D_{R}$ 上的平凡 $G$-主丛。


设 $T\subset G$ 是分裂极大环面，Weyl 群为 $W$。回忆：仿射 Grassmann 簇 $\Gr_{G}$ 上的 $L^{+}G$-轨道由余权在 Weyl 群作用下的轨道 $\xcoch(T)/W$ 标号；对 $\l\in \xcoch(T)$，我们指定包含 $t^{\l}\in LT$ 的那条 $L^{+}G$-轨道，它只依赖于 $\l$ 的 $W$-轨道。

选取包含 $T$ 的 Borel 子群，它定义了支配余权 $\xcoch(T)^{+}$。于是有典范双射 $\xcoch(T)^{+}\isom \xcoch(T)/W$，这使我们可以用 $\xcoch(T)^{+}$ 中的元素唯一地标号 $\Gr_{G}$ 上的任一条 $L^{+}G$-轨道。对 $\l\in \xcoch(T)^{+}$，令 $\Gr_{G, \l}$ 为 $t^{\l}$ 所在的 $L^{+}G$-轨道，视作局部闭的约化子概形。我们有 $\dim\Gr_{G,\l}=\j{2\r,\l}$，其中 $2\r$ 是正根之和。

在 $\xcoch(T)^{+}$ 上定义偏序 $\le$ 如下：$\l'\le \l$ 当且仅当 $\l-\l'$ 是关于 $B$ 的正余根之和。已知 $\Gr_{G,\l}$ 的闭包为 $\Gr_{G,\le \l}=\cup_{\l'\in \xcoch(T)^{+}, \l'\le \l}\Gr_{G,\l'}$。


对闭点 $x\in |X|$，类似地定义 $\Gr_{G,x}$ 为表示 $R\mapsto \Gr_{G,x}(R)=(L_{x}G)(R)/(L^{+}_{x}G)(R)$ 的 fpqc 层化的 ind-概形。选定 $x$ 处的单值化子 $\vp_{x}$ 后，得到同构 $\Gr_{G,x}\cong \Gr_{G}\ot_{k}k_{x}$。当 $x$ 变动时，$\Gr_{G,x}$ 构成 $X$ 上的一个族，这是 Beilinson–Drinfeld Grassmann 簇的一个特例，下面就来回忆它。

\begin{exam} 当 $G=\GL_{n}$ 时，$\Gr_{G}(k)$ 是 $k\lr{t}^{\op n}$ 中秩 $n$ 的自由 $k\tl{t}$-子模的集合。设整数 $d_{1}\ge d_{2}\ge\cdots \ge d_{n}$，令 $\l=(d_{1},d_{2},\cdots, d_{n})$ 并视作对角环面 $T\subset G$ 的支配余权，则点 $t^{\l}$ 对应 $k\tl{t}$-子模
\begin{equation*}
t^{d_{1}}k\tl{t}\op t^{d_{2}}k\tl{t}\op\cdots\op t^{d_{n}}k\tl{t}\subset k\lr{t}^{\op n}.
\end{equation*}

当 $d_{n}\ge0$ 时，$\Gr_{G,\l}(k)$ 是 $k\tl{t}$-子模 $\L\subset k\tl{t}^{\op n}$ 的集合，使得作为 $k\tl{t}$-模有 $k\tl{t}^{\op n}/\L\cong \op_{i=1}^{n}k\tl{t}/(t^{d_{i}})$。

当 $\l=(d,0,\cdots, 0)$ 且 $d\ge0$ 时，$\Gr_{G,\le\l}(k)$ 是使 $\dim_{k}k\tl{t}^{\op n}/\L=d$（而非 $\le d$）的 $k\tl{t}$-子模 $\L\subset k\tl{t}^{\op n}$ 的集合。
\end{exam}



\ssszh{Beilinson–Drinfeld Grassmann 簇}\label{sss:BD}
令 $\cE_{\triv}\in \Bun_{G}(k)$ 为 $X$ 上的平凡 $G$-主丛。用笛卡儿方图
\begin{equation*}
\begin{tikzcd}
\Gr^{(I_{1},\cdots, I_{r})}_{G} \arrow[r] \arrow[d] & \Hk^{(I_{1},\cdots, I_{r})}_{G} \arrow[d, "p_{r}"] \\
\Spec k \arrow[r, "\cE_{\triv}"] & \Bun_{G}
\end{tikzcd}
\end{equation*}
定义 {\em Beilinson–Drinfeld Grassmann 簇} $\Gr^{(I_{1},\cdots, I_{r})}_{G}$。于是有腿映射 $\Gr^{(I_{1},\cdots, I_{r})}_{G}\to X^{I}$。当 $I$ 是单元集 $\{1\}$ 时，$\Gr^{\{1\}}_{G}\to X$ 在闭点 $x\in |X|$ 处的纤维典范同构于 $\Gr_{G,x}$。

\ssszh{相对位置映射}
可以构造商叠 $\QGr_{G}:=L^{+}G\bs \Gr_{G}$，它是 $\Hk^{\{1\}}_{G}$ 的局部版本：对 $k$-代数 $R$，$\QGr_{G}$ 的 $R$-点是三元组 $(\cE_{0}, \cE_{1}, \a)$，其中 $\cE_{i}$ 是 $D_{R}=\Spec R\tl{t}$ 上的 $G$-主丛，$\a: \cE_{0}|_{D^{\times}_{R}}\isom \cE_{1}|_{D^{\times}_{R}}$。于是 $\QGr_{G}$ 的点与 $\xcoch(T)^{+}$ 一一对应。对 $\l\in \xcoch(T)^{+}$，令 $\QGr_{G, \le \l}$ 为 $\Gr_{G,\le \l}$ 的像，它是 $\QGr_{G}$ 的闭约化子叠。

令 $\Aut(D)$ 为 $k\tl{t}$ 的自同构构成的 $k$ 上 pro-代数群。则 $\Aut(D)$ 通过坐标变换右作用在 $\QGr_{G}$ 上：$\th: D\to D$ 把 $(\cE_{0},\cE_{1}, \a)$ 送到 $(\th^{*}\cE_{0}, \th^{*}\cE_{1}, \th^{*}\a)$。显然 $\QGr_{G, \le\l}$ 在 $\Aut(D)$ 的作用下不变。


对 Hecke 叠 $\Hk^{\{1\}}_{G,\bK}$，如下定义求值映射
\begin{equation*}
\ev^{\{1\}}_{\Hk}: \Hk^{\{1\}}_{G,\bK}\to \QGr_{G}/\Aut(D)
\end{equation*}
只需构造一个 $\Aut(D)$-主丛 $\wt\Hk^{\{1\}}_{G,\bK}\to \Hk^{\{1\}}_{G,\bK}$，并配以 $\Aut(D)$-等变映射 $\wt\ev^{\{1\}}_{\Hk}: \wt\Hk^{\{1\}}_{G,\bK}\to \QGr_{G}$。定义 $\wt\Hk^{\{1\}}_{G,\bK}$ 为如下模叠：其 $R$-点为 $(x, t_{x}, \cE_{0}, \cE_{1}, \a)$，其中 $(x,\cE_{0},\cE_{1}, \a)\in\Hk^{\{1\}}_{G,\bK}(R)$，而 $t_{x}$ 是 $X\ot R$ 沿 $\G(x)$ 的形式完备化上的正则函数，且对 $s\in \Spec R$ 的每个几何点，$t_{x}$ 都是 $\wh\cO_{X\times s, x_{s}}$ 的单值化子。于是 $t_{x}$ 给出连续同构 $\wh \cO_{X\ot R, \G(x)}\cong R\tl{t}$，从而给出映射 $t^{\sh}_{x}: \Spec R\tl{t}\to X\ot R$。映射 $\wt\ev^{\{1\}}_{\Hk}$ 便把 $(x, t_{x}, \cE_{0}, \cE_{1}, \a)$ 送到 $(t^{\sh, *}_{x}\cE_{0}, t^{\sh, *}_{x}\cE_{1}, t^{\sh, *}_{x}\a)$。

\ssszh{有界的 Hecke 叠}\label{sss:bd Hk}
对 $\l\in \xcoch(T)^{+}$，定义闭子叠
\begin{equation*}
\Hk^{\{1\}, \le \l}_{G,\bK}\subset \Hk^{\{1\}}_{G,\bK}
\end{equation*}
为 $\QGr_{G,\le \l}/\Aut(D)\subset \QGr_{G}/\Aut(D)$ 在上述映射 $\ev^{\{1\}}_{\Hk}$ 下的原像。粗略地说，$\Hk^{\{1\}}_{G,\bK}$ 的 $S$-点 $(x,\cE_{0},\cE_{1}, \a)$ 落在 $\Hk^{\{1\}, \le \l}_{G,\bK}$ 中，当且仅当 $\a$ 的相对位置 $\le \l$。

现在设有限集 $I$ 有分解 $I=I_{1}\coprod I_{2}\coprod \cdots\coprod I_{r}$。令 $\un\l=(\l_{i})_{i\in I}\in (\xcoch(T)^{+})^{I}$。我们想定义 $\Hk^{(I_{1},\cdots, I_{r})}_{G,\bK}$ 的闭子叠 $\Hk^{(I_{1},\cdots, I_{r}), \le \un\l}_{G,\bK}$。构造分几步完成：
\begin{enumerate}
\item 令 $U^{I_{j}}_{\disj}\subset U^{I_{j}}$ 为使 $\G(x_{i})$ 两两不相交的 $(x_{i})_{i\in I_{j}}$ 构成的开子集。则有映射
\begin{equation}\label{ev disj}
\ev^{(I_{1},\cdots, I_{r})}_{\Hk,\disj}: \Hk^{(I_{1},\cdots, I_{r})}_{G,\bK}|_{U^{I_{1}}_{\disj}\times\cdots\times U^{I_{r}}_{\disj}}\to (\QGr_{G}/\Aut(D))^{I}
\end{equation}
其第 $i$ 个分量记录 $\a_{j}$（其中 $i\in I_{j}$）沿 $\G(x_{i})$ 的相对位置。定义
\begin{equation*}
\Hk^{(I_{1},\cdots, I_{r}), \le \un\l}_{G,\bK, \disj}\subset \Hk^{(I_{1},\cdots, I_{r})}_{G,\bK}|_{U^{I_{1}}_{\disj}\times\cdots\times U^{I_{r}}_{\disj}}
\end{equation*}
为 $\prod_{i\in I}\QGr_{G,\le \l_{i}}/\Aut(D)$ 在 $\ev^{(I_{1},\cdots, I_{r})}_{\Hk,\disj}$ 下的原像。


\item 接着定义 $\Hk^{(I_{1},\cdots, I_{r}), \le \un\l}_{G,\bK}$ 为 $\Hk^{(I_{1},\cdots, I_{r}), \le \un\l}_{G,\bK, \disj}$ 在 $\Hk^{(I_{1},\cdots, I_{r})}_{G,\bK}$ 中的闭包。
\end{enumerate}

定义 Beilinson–Drinfeld Grassmann 簇的有界版本 $\Gr^{(I_{1},\cdots, I_{r}), \le\un\l}_{G}$ 为 $\Hk^{(I_{1},\cdots, I_{r}), \le\un\l}_{G}$ 在典范映射 $\Gr^{(I_{1},\cdots, I_{r})}_{G}\to \Hk^{(I_{1},\cdots, I_{r})}_{G}$ 下的原像。

\begin{exam} 设 $G=\GL_{n}$，并对所有 $i\in I$ 取 $\l_{i}=(1,0,\cdots, 0)$。
\begin{enumerate}
\item 迭代版本。设 $I=\{1,\cdots, r\}$ 分解为单点集 $I_{i}=\{i\}$。则 $\Hk^{(\{1\},\cdots,\{ r\}),\le\un\l}_{G}$ 的 $S$-点是
$$((x_{i})_{1\le i\le r}, (\cV_{j})_{0\le j\le r}, (\a_{j})_{1\le j\le r}),$$
其中 $\cV_{j}$ 是 $X\times S$ 上秩 $n$ 的向量丛，$\a_{j}: \cV_{j-1}\incl \cV_{j}$ 是凝聚层的单射，其余核在图像 $\G(x_{j})$ 上局部自由且秩为 $1$，$1\le j\le r$。

\item 非迭代版本。$\Hk^{I,\le\un\l}_{G}$ 的 $S$-点是 $((x_{i})_{i\in I}, \cV_{0}, \cV_{1}, \a)$，其中 $\cV_{0}$ 与 $\cV_{1}$ 是 $X\times S$ 上秩 $n$ 的向量丛，$\a: \cV_{0}\incl \cV_{1}$ 是凝聚层的单射，使得 $\det(\a): \det(\cV_{0})\to \det(\cV_{1})$ 的除子作为 $X\times S$ 上的除子是 $\sum_{i\in I}\G(x_{i})$。
\end{enumerate}

\end{exam}


\begin{rmk} 我们的记号与 \cite[\S1]{VLaff} 不同。Hecke 叠 $\Hk^{(I_{1},\cdots, I_{r})}_{G,\bK}$ 在 {\em 同上} 中记作 $\textup{Hecke}_{\bK,I}^{(I_{1},\cdots, I_{r})}$。在 \cite[定义 1.2]{VLaff} 中引入了迭代 Hecke 叠的两个有界版本，分别记作 $\textup{Hecke}^{(I_{1},\cdots, I_{r})}_{\bK, I, \un\l}$ 与 $\textup{Hecke}^{(I_{1},\cdots, I_{r})}_{\bK, I, \lesssim\un\l}$。可以证明，我们的 Hecke 叠 $\Hk^{(I_{1},\cdots, I_{r}),\le\un\l}_{G, \bK}$ 是 $\textup{Hecke}^{(I_{1},\cdots, I_{r})}_{\bK, I, \un\l}$ 的约化结构，而后者又包含于 $\textup{Hecke}^{(I_{1},\cdots, I_{r})}_{\bK, I, \lesssim\un\l}$。
\end{rmk}

%%%%%%%%%%%%%%
\subsection{$G$-shtuka}\label{ss:Sht}
对 $k=\FF_{q}$ 上任意概形 $Y$，用 $\Fr_{Y}$ 或简写 $\Fr: Y\to Y$ 表示把正则函数取 $q$ 次幂的 Frobenius 态射。

仍固定有限个位点组成的集合 $\Sigma\subset |X|$ 以及水平群 $\bK=(\bK_{x})_{x\in \Sig}$。令 $U=X-\Sigma$。

\ssszh{$\Sht_{G}$ 的定义}
设有限集 $I$ 有分解 $I=I_{1}\coprod \cdots\coprod I_{r}$。用笛卡儿方图
\begin{equation}\label{defn Sht}
\begin{tikzcd}
\Sht_{G,\bK}^{(I_{1},\cdots, I_{r})} \arrow[r] \arrow[d] & \Hk_{G,\bK}^{(I_{1},\cdots, I_{r})} \arrow[d, "{({p_0},{p_r})}"'] \\
\Bun_{G,\bK} \arrow[r, "{(\id,\Fr)}"] & \Bun_{G,\bK}\times \Bun_{G,\bK}
\end{tikzcd}
\end{equation}
定义 Artin 叠 $\Sht_{G,\bK}^{(I_{1},\cdots, I_{r})}$。

由构造，有记录腿的映射
\begin{equation*}
\pi:=\pi^{(I_{1},\cdots, I_{r})}_{G,\bK}: \Sht_{G,\bK}^{(I_{1},\cdots, I_{r})}\to U^{I}.
\end{equation*}


\ssszh{点函子}\label{sss:Sht points}
按定义，对 $k$ 上任意测试概形 $S$，$\Sht_{G,\bK}^{I}$ 的 $S$-点是多元组
\begin{equation*}
\xi=((x_{i})_{i\in I}, \cE, \a: \cE\dr {}^{\t}\cE)
\end{equation*}
的群胚，其中
\begin{itemize}
\item 对 $i\in I$，$x_{i}\in U(S)$，其图像为 $\G(x_{i})\subset U\times S$。这些 $x_{i}$ 称为 shtuka $\xi$ 的 {\em 腿}。
\item $\cE$ 是 $X\times S$ 上的 $G$-主丛，并沿 $\{x\}\times S$ 带有 $\bK_{x}$-水平结构。
\item ${}^{\t}\cE=(\id_{X}\times \Fr_{S})^{*}\cE$。
\item $\a$ 是保持 $\bK$-水平结构的 $G$-主丛同构 $\cE|_{X\times S-\cup_{i\in I}\G(x_{i})}\isom ({}^{\t}\cE)|_{X\times S-\cup_{i\in I}\G(x_{i})}$。
\end{itemize}


更一般地，对分解 $I=I_{1}\coprod \cdots\coprod I_{r}$，$\Sht_{G,\bK}^{(I_{1},\cdots, I_{r})}$ 的 $S$-点是多元组
\begin{equation*}
((x_{i})_{i\in I},(\cE_{j})_{0\le j\le r}, (\a_{j}: \cE_{j-1}\dr \cE_{j})_{1\le j\le r}, \io)
\end{equation*}
的群胚，其中
\begin{itemize}
\item 对 $i\in I$，$x_{i}\in U(S)$，其图像为 $\G(x_{i})\subset X\times S$。
\item 对 $0\le j\le r$，$\cE_{j}$ 是 $X\times S$ 上的 $G$-主丛，并沿 $\{x\}\times S$ 带有 $\bK_{x}$-水平结构。
\item 对 $1\le j\le r$，$\a_{j}$ 是保持 $\bK$-水平结构的 $G$-主丛同构
\begin{equation*}
\cE_{j-1}|_{X\times S-\cup_{i\in I_{j}}\G(x_{i})}\isom \cE_{j}|_{X\times S-\cup_{i\in I_{j}}\G(x_{i})}.
\end{equation*}
\item $\io$ 是保持 $\bK$-水平结构的 $G$-主丛同构
\begin{equation*}
\io:\cE_{r}\isom {}^{\t}\cE_{0}=(\id_{X}\times \Fr_{S})^{*}\cE_{0}.
\end{equation*}
\end{itemize}

\begin{exam}\label{ex:Sht no leg}
当 $I=\vn$ 时，按定义 $\Sht^{\vn}_{G,\bK}$ 是对角态射 $\D: \Bun_{G,\bK}\to \Bun_{G,\bK}^{2}$ 与 Frobenius 图像 $(\id,\Fr): \Bun_{G,\bK}\to \Bun_{G,\bK}^{2}$ 的纤维积。一般事实是：对 $k$ 上任意 Artin 叠 $\frX$，$\D: \frX\to \frX^{2}$ 与 Frobenius 图像 $(\id,\Fr):\frX\to \frX^{2}$ 的纤维积同构于群胚 $\frX(k)$，视作 $k$ 上的离散叠。在我们的情形有
\begin{equation*}
\Sht^{\vn}_{G,\bK}=\coprod_{\cE\in \Bun_{G,\bK}(k)}[(\Spec k)/\Aut(\cE)(k)].
\end{equation*}
\end{exam}



\ssszh{有界版本}
令 $\un\l=(\l_{i})_{i\in I}\in (\xcoch(T)^{+})^{I}$。用笛卡儿方图
\begin{equation*}
\begin{tikzcd}
\Sht_{G,\bK}^{(I_{1},\cdots, I_{r}),\le\un\l}\arrow[r]\arrow[d] & \Hk_{G,\bK}^{(I_{1},\cdots, I_{r}),\le\un\l}\arrow[d, "{(p_{0},p_{r})}"] \\
\Bun_{G,\bK}\arrow[r, "{(\id,\Fr)}"] & \Bun_{G,\bK}\times \Bun_{G,\bK}
\end{tikzcd}
\end{equation*}
定义 $\Sht^{(I_{1},\cdots, I_{r})}_{G,\bK}$ 的有界版本 $\Sht^{(I_{1},\cdots, I_{r}), \le \un\l}_{G,\bK}$。于是 $\Sht^{(I_{1},\cdots, I_{r}), \le \un\l}_{G,\bK}$ 是 $\Sht_{G,\bK}^{(I_{1},\cdots, I_{r})}$ 的闭子叠。



\begin{rmk} 我们的记号与 \cite[定义 2.1]{VLaff} 不同。叠 $\Sht^{(I_{1},\cdots, I_{r})}_{G,\bK}$ 在 {\em 同上} 中记作 $\textup{Cht}_{\bK,I}^{(I_{1},\cdots, I_{r})}$。至于有界版本，叠 $\textup{Cht}_{\bK,I, \un\l}^{(I_{1},\cdots, I_{r})}$ 是在 {\em 同上} 中引入的，并且存在闭嵌入 $\Sht^{(I_{1},\cdots, I_{r}), \le \un\l}_{G,\bK}\incl \textup{Cht}_{\bK,I, \un\l}^{(I_{1},\cdots, I_{r})}$，它在取约化结构后诱导同构。
\end{rmk}


\begin{exam}\label{ex:Gm} 设 $G=\GL_{1}=\Gm$，$I$ 为有限集。此时一个界由整数列 $\un\l:=(\l_{i})_{i\in I}\in \ZZ^{I}$ 给出。注意 $\ZZ^{I}$ 上的偏序是平凡的：$\l'\le \l$ 当且仅当 $\l'=\l$。叠 $\Sht^{I, \le \un\l}_{\Gm}$ 另有如下描述：它适合笛卡儿方图
\begin{equation*}
\begin{tikzcd}
\Sht^{I, \le \un\l}_{\Gm} \arrow[r]\arrow[d] & \Pic_{X}\arrow[d, "\bL_{\Pic_{X}}"] \\
X^{I}\arrow[r, "\AJ_{\un\l}"] & \Pic_{X}
\end{tikzcd}
\end{equation*}
这里 $\bL_{\Pic_{X}}$ 是 $\Pic_{X}$ 的 Lang 同源，它把线丛 $\cL$ 送到 ${}^{\t}\cL\ot \cL^{-1}$。Abel–Jacobi 映射 $\AJ_{\un\l}$ 把 $(x_{i})_{i\in I}$ 送到线丛 $\cO_{X}(\sum_{i\in I}\l_{i}x_{i})$。

由于 $\bL_{\Pic_{X}}$ 的像落在 $\Pic^{0}_{X}$ 中，而 $\AJ_{\un\l}$ 落在 $\Pic^{|\un\l|}_{X}$ 中（其中 $|\un\l|=\sum_{i\in I}\l_{i}$），可见当 $|\un\l|\ne0$ 时 $\Sht^{I, \le \un\l}_{G} =\vn$。
\end{exam}

\begin{exam}\label{ex:Dr} 设 $G=\GL_{n}$，并设 $D=\sum_{x\in \Sig}m_{x}x$ 是 $X$ 上的有效除子。令 $\bK_{D}$ 为 $x\in \Sig$ 处主同余子群 $\bK_{x,m_{x}}$ 构成的族。令 $I=\{1,2\}$，并取分解 $I_{1}=\{1\}$、$I_{2}=\{2\}$。令 $\l_{1}=(0,\cdots, 0,-1)$ 与 $\l_{2}=(1,0,\cdots,0)$ 为两个极小余权，$\un\l=(\l_{1},\l_{2})$。则 $\Sht^{(I_{1},I_{2}), \le\un\l}_{G,\bK_{D}}(S)$ 分类如下数据
\begin{equation*}
\begin{tikzcd}
\cV_{0} & \cV_{1}\arrow[l, "\a_{1}"'] \arrow[r, "\a_{2}"] & \cV_{2}\arrow[r, "\io"{above}, "{\sim}"{below}] & {}^{\t}\cV_{0}
\end{tikzcd}
\end{equation*}
其中 $\cV_{i}$ 是 $X\times S$ 上秩 $n$ 的向量丛，并沿 $D\times S$ 带有平凡化；$\a_{1}$ 与 $\a_{2}$ 是与 $D\times S$ 上的平凡化相容的凝聚层单射，其余核分别在映射 $x_{1}: S\to U=X-\Sig$ 与 $x_{2}: S\to U$ 的图像上局部自由。这正是 Drinfeld 在 \cite{Dr-FBun} 中最初考虑的那类模叠。
\end{exam}

\ssszh{容许性}
令 $\Om=\xcoch(T)/\ZZ\Phi^{\vee}$，其中 $\ZZ\Phi^{\vee}$ 是余根格。则 $\Om$ 是 $\Gr_{G}$（以及 $\QGr_{G}$）的连通分支群。由 \cite[引理 2.2]{Var}，连通分支集 $\pi_{0}(\Bun_{G})$ 也可典范地等同于 $\Om$。更一般地，若每个 $\bK_{x}$ 都包含在 $L_{x}G$ 的某个抛孔子群中，则 $\pi_{0}(\Bun_{G,\bK})$ 仍可典范地等同于 $\Om$，其中平凡丛所在的分支对应 $0\in \Om$。


令 $\un\l=(\l_{i})_{i\in I}\in (\xcoch(T)^{+})^{I}$，$\om$ 为 $\sum_{i\in I}\l_{i}\in \xcoch(T)$ 在 $\Om$ 中的像。若 $\om=0\in \Om$，则称 $\un\l$ 是 {\em $G$-容许的}（见 \cite[定义 2.14]{Var}）。

\begin{lem}[Varshavsky {\cite[命题 2.16(d)]{Var}}] 设对每个 $x\in \Sig$，$\bK_{x}$ 都包含在某个抛孔子群中。设 $I=I_{1}\coprod\cdots\coprod I_{r}$ 是有限集 $I$ 的分解，$\un\l=(\l_{i})_{i\in I}\in (\xcoch(T)^{+})^{I}$。

则 $\Sht^{(I_{1},\cdots, I_{r}), \le\un\l}_{G,\bK}$ 非空当且仅当 $\un\l$ 是 $G$-容许的。
\end{lem}
\begin{proof}[证明概要]

对 $\g\in \Om\cong \pi_{0}(\Bun_{G,\bK})$，以 $\Bun_{G, \bK}^{\g}$ 记 $\Bun_{G,\bK}$ 相应的连通分支。把映射 $(p_{0},p_{r}): \Hk^{(I_{1},\cdots, I_{r})}_{G,\bK}\to \Bun_{G,\bK}\times\Bun_{G,\bK}$ 限制到 $\Hk^{(I_{1},\cdots, I_{r}), \le\un\l}_{G,\bK}$ 上，得
\begin{equation*}
(p^{\le\un\l}_{0},p^{\le\un\l}_{r}):  \Hk^{(I_{1},\cdots, I_{r}), \le\un\l}_{G,\bK}\to \coprod_{\g\in \Om} \Bun_{G,\bK}^{\g}\times \Bun^{\g+\om}_{G,\bK}
\end{equation*}
其中 $\om$ 是 $\sum\l_{i}$ 在 $\Om$ 中的像。

若 $\om\ne0$，则 $(p^{\le\un\l}_{0},p^{\le\un\l}_{r})$ 的像与 $\Bun_{G,\bK}^{\g}\times \Bun^{\g}_{G,\bK}$（对任意 $\g\in \Om$）都不相交。然而 $\Bun_{G,\bK}$ 上 Frobenius 的图像落在 $\coprod_{\g\in \Om} \Bun_{G,\bK}^{\g}\times \Bun^{\g}_{G,\bK}$ 之上。由笛卡儿方图 \eqref{defn Sht} 即知 $\Sht^{(I_{1},\cdots, I_{r}), \le\un\l}_{G,\bK}=\vn$。

反之，设 $\un\l$ 是 $G$-容许的。由于 $\Hk^{(I_{1},\cdots, I_{r}),\le\un\l}_{G, \bK}\to \Hk^{I,\le\un\l}_{G,\bK}$ 满射，映射 $\Sht^{(I_{1},\cdots, I_{r}),\le\un\l}_{G, \bK}\to \Sht^{I,\le\un\l}_{G,\bK}$ 也满射。因此只需证明 $\Sht^{I,\le\un\l}_{G,\bK}\ne\vn$。

令 $\l=\sum_{i}\l_{i}$。因 $\un\l$ 容许，支配余权 $\l$ 落在 $G$ 的余根格中，故 $0\le \l$。把 $\Sht^{I,\le\un\l}_{G,\bK}$ 限制到对角线 $\D(U)\subset U^{I}$，得 $\Sht^{I,\le\un\l}_{G,\bK}|_{\D(U)}\cong \Sht^{\{1\},\le\l}_{G,\bK}$。这一点将在 \S\ref{ss:geom fact} 中讨论。于是只需证明 $\Sht^{\{1\},\le\l}_{G,\bK}\ne\vn$。因 $0\le \l$，有 $\Sht^{\{1\},\le0}_{G,\bK}\subset \Sht^{\{1\},\le\l}_{G,\bK}$。显然 $\Sht^{\{1\},\le0}_{G,\bK}=\Sht^{\vn}_{G,\bK}\times U$，由例 \ref{ex:Sht no leg} 知它非空。
\end{proof}
